\chapter[Arquitetura de Hardware]{Arquitetura de Hardware}
\label{cap:hardware}

Este capítulo descreve os componentes físicos do robô, os critérios que
orientaram a escolha de cada um, o princípio de funcionamento de cada
componente e a forma como eles se integram. O robô é uma plataforma de
tração nas quatro rodas (4WD) que deve mapear um ambiente com um sensor
LiDAR e uma câmera de profundidade e, em seguida, deslocar-se de forma
autônoma até um ponto escolhido no mapa. Esses dois objetivos determinam
as exigências sobre o hardware: o robô precisa medir o próprio
deslocamento (odometria), perceber o ambiente ao redor e transmitir esses
dados a um computador com capacidade para processá-los.

A Tabela~\ref{tab:componentes} resume os componentes adotados. As seções
seguintes detalham cada um deles.

\begin{table}[htb]
	\centering
	\caption{Componentes de hardware do robô}
	\label{tab:componentes}
	\begin{tabular}{|l|l|}
		\hline
		\textbf{Função} & \textbf{Componente} \\ \hline
		Processamento embarcado & Raspberry Pi 3 (1~GB de RAM) \\ \hline
		Estrutura & Chassi 4WD de 20~cm $\times$ 21~cm \\ \hline
		Atuação & 4 motores JGB37-520 (12~V, 333~rpm) \\ \hline
		Odometria & Encoders Hall em quadratura (11~PPR) \\ \hline
		Acionamento & 2 drivers Cytron MD10C R3 \\ \hline
		Orientação & IMU Bosch BNO055 (UART) \\ \hline
		Percepção 2D & LiDAR Slamtec RPLIDAR C1 \\ \hline
		Percepção 3D & Câmera Orbbec Astra Pro \\ \hline
		Comunicação & Adaptador Wi-Fi USB (RTL8188EU) \\ \hline
		Energia & Bateria LiPo 4S com BMS e dois conversores \textit{buck} \\ \hline
	\end{tabular}
	\legend{Fonte: Autor.}
\end{table}

\section{Unidade de processamento embarcada}

A unidade de processamento embarcada é um Raspberry Pi 3, um computador de
placa única com processador ARM Cortex-A53 de quatro núcleos e 1~GB de
memória RAM \cite{raspberrypi}. Ele foi escolhido por rodar um sistema
Linux completo, o que permite usar o ROS~2 e as bibliotecas dos sensores
sem precisar escrever drivers do zero. Ele também tem uma barra de 40
pinos de propósito geral (GPIO), com dois canais de PWM por hardware, uma
interface UART e portas USB. Essas interfaces bastam para ligar todos os
componentes do robô sem placas adicionais.

O Raspberry Pi 3 tem pouca capacidade de processamento e de memória para
executar ao mesmo tempo o mapeamento simultâneo (SLAM), a navegação e a
visualização. Por isso, ele foi usado apenas como computador de interface
com o hardware. Ele lê os sensores, aciona os motores e envia os dados por
Wi-Fi a um notebook, onde roda o processamento pesado. Essa divisão está
descrita no Capítulo~\ref{cap:software}. Para liberar memória, a placa
opera sem interface gráfica e é acessada remotamente por SSH.

\section{Chassi, motores e encoders}

A plataforma mecânica é um chassi 4WD de 20~cm de largura por 21~cm de
comprimento, com rodas de 8~cm de diâmetro. O robô é do tipo
\textit{skid-steer}: as rodas não giram em torno de um eixo vertical, e as
curvas são feitas girando os dois lados com velocidades diferentes. Essa
configuração foi adotada por ser mecanicamente simples e por ter boa tração
em pisos irregulares. Em troca, as rodas escorregam lateralmente durante
as curvas, e isso afeta a odometria, como discutido na
Seção~\ref{sec:integracao-sw}.

Cada roda é movida por um motorredutor JGB37-520 de 12~V com velocidade
nominal de 333~rpm no eixo de saída. A escolha considerou a tensão,
compatível com a bateria por meio de um conversor, o torque oferecido pela
caixa de redução e, principalmente, o encoder já integrado ao motor. Com
rodas de 8~cm, a velocidade linear teórica sem carga é
\begin{equation}
	\label{eq:vmax}
	v_{max} = \pi D \frac{n}{60} = \pi \cdot 0{,}08 \cdot \frac{333}{60}
	\approx 1{,}4~\text{m/s},
\end{equation}
em que $D$ é o diâmetro da roda em metros e $n$ é a velocidade do eixo em
rpm. Esse valor está bem acima do necessário para mapear ambientes
internos, o que deixa margem para operar em velocidades baixas, que
favorecem a qualidade do mapa.

O encoder de cada motor tem dois sensores de efeito Hall dispostos em
torno de um ímã preso ao eixo do motor, antes da redução. Quando o ímã
gira, cada sensor gera uma onda quadrada com 11 pulsos por volta do eixo
do motor (PPR). Os dois sinais, chamados A e B, ficam defasados de
90\textdegree. Esse arranjo é chamado de quadratura. A ordem em que os
sinais mudam de nível indica o sentido de giro, e contar as quatro bordas
de cada ciclo (decodificação $\times$4) multiplica a resolução por quatro.
O número de contagens por volta da roda é então
\begin{equation}
	\label{eq:cpr}
	N_{CPR} = 4 \cdot PPR \cdot i,
\end{equation}
em que $i$ é a relação de redução da caixa. Com esse valor, a distância
percorrida pela roda em um intervalo com $\Delta N$ contagens é
\begin{equation}
	\label{eq:deslocamento}
	\Delta s = \pi D \frac{\Delta N}{N_{CPR}}.
\end{equation}
% TODO: preencher com o CPR_SAIDA medido pelo teste_motores.py encoder.
O valor de $N_{CPR}$ é medido experimentalmente girando a roda um número
conhecido de voltas, conforme descrito na Seção~\ref{sec:bancada}.

Os dois motores de cada lado recebem o mesmo comando, então medir uma
roda por lado basta para estimar a velocidade daquele lado. Por isso, só
os encoders dos motores dianteiros estão ligados ao Raspberry Pi. Os
traseiros recebem apenas alimentação. Assim são usadas quatro entradas
digitais em vez de oito, e a carga de interrupções no processador também
cai pela metade.

\section{Drivers de motor}

Os pinos do Raspberry Pi fornecem poucos miliampères a 3,3~V. Por isso, os
motores são acionados por dois drivers Cytron MD10C R3, um para cada lado
do robô, com os dois motores daquele lado ligados em paralelo. O MD10C é
uma ponte H, um circuito com quatro chaves eletrônicas (MOSFETs) que
liga os terminais do motor à alimentação em uma polaridade ou na outra
\cite{cytronmd10c}. Ele foi escolhido por operar na faixa de tensão do
robô, por suportar com folga a corrente de dois motores em paralelo e por
ser controlado com apenas dois sinais lógicos compatíveis com 3,3~V.
% TODO: conferir no manual do MD10C R3 a faixa de tensão (5-30 V) e a
% corrente continua e de pico, e citar os valores aqui.

O controle segue o esquema sinal-magnitude. O pino DIR define o sentido
de giro, e o pino PWM recebe uma onda quadrada cujo ciclo de trabalho
(\textit{duty cycle}) define a fração do tempo em que a ponte fica
ligada. Como o motor responde à média da tensão aplicada, o ciclo de
trabalho controla a velocidade. O PWM é gerado pelos periféricos de PWM
por hardware do Raspberry Pi nos pinos GPIO12 e GPIO13, a 10~kHz. A
geração por hardware não depende do agendamento do sistema operacional,
então a frequência fica estável mesmo com o processador ocupado. Nos
pinos de PWM há resistores de 10~k$\Omega$ para o terra, que mantêm os
motores parados durante a inicialização do sistema, enquanto os pinos
ainda não foram configurados.

\section{Unidade de medição inercial}

A odometria por encoders é pouco confiável para estimar a orientação de
um robô \textit{skid-steer}, porque as rodas escorregam nas curvas. Para
medir a rotação de forma independente, foi usada a unidade de medição
inercial (IMU) Bosch BNO055. Ela reúne acelerômetro, giroscópio e
magnetômetro de três eixos e tem um microcontrolador interno que executa
a fusão dos sensores e fornece a orientação já calculada, na forma de
quatérnios ou ângulos de Euler \cite{bno055}. Isso tira do Raspberry Pi o
trabalho de calibração e de fusão e é o principal motivo da escolha.

O BNO055 aceita I$^2$C ou UART, e a interface é escolhida pelos pinos PS0
e PS1. Foi adotada a UART (PS0 em nível baixo e PS1 em nível alto),
ligada aos pinos GPIO14 e GPIO15 com TX e RX cruzados. O sensor usa
\textit{clock stretching} no barramento I$^2$C, recurso que o controlador
I$^2$C do Raspberry Pi não trata de forma confiável. A UART evita esse
problema. O sensor é alimentado pelo pino de 3,3~V da própria placa.

\section{Sensor LiDAR}

O sensor principal de mapeamento é o LiDAR Slamtec RPLIDAR C1. Ele gira
um feixe laser em 360\textdegree{} e mede a distância até o obstáculo mais
próximo em cada direção pelo tempo de voo da luz \cite{rplidarc1}. No
robô, o driver identificou o sensor operando em modo padrão com 5~mil
medidas por segundo e 10 voltas por segundo, ou seja, cerca de 500 pontos
por volta. Cada volta forma uma varredura 2D do plano em que o sensor
está, a 10~cm do chão, e essa varredura é a entrada do algoritmo de SLAM.
% TODO: citar o alcance e a resolucao angular do datasheet do C1.

O C1 foi escolhido por ter custo baixo, ter driver oficial para ROS~2
(\texttt{sllidar\_ros2}) e alcance compatível com ambientes internos. Ele
se comunica por UART a 460800~baud por meio de um adaptador USB-serial
CP210x, ligado a uma porta USB do Raspberry Pi. O sensor foi instalado
aproximadamente no centro do chassi, de modo que a posição medida fique
próxima do centro de rotação do robô.

\section{Câmera de profundidade}

Para perceber obstáculos fora do plano do LiDAR, como mesas, degraus e
objetos baixos, o robô tem uma câmera de profundidade Orbbec Astra Pro.
Ela projeta um padrão de luz infravermelha estruturada sobre a cena e
calcula a distância de cada pixel pela deformação desse padrão, observada
por uma câmera infravermelha \cite{orbbecastra}. A câmera também tem um
sensor RGB convencional. No robô, as duas partes aparecem como
dispositivos USB distintos: o RGB pela classe UVC e a profundidade pelo
SDK OpenNI2 da fabricante.

Nos testes, a profundidade foi lida a cerca de 30 quadros por segundo.
Abaixo de aproximadamente 0,6~m a câmera não retorna medida, e esse é o
seu alcance mínimo. A câmera foi escolhida pela disponibilidade e por ter
drivers para ROS~2. O uso dela na navegação está previsto para uma etapa
posterior, com resolução reduzida (320$\times$240 a 15 quadros por
segundo) para caber na capacidade do Raspberry Pi 3 e na banda da rede.

\section{Alimentação}

A energia vem de uma bateria LiPo de quatro células em série (4S), com
tensão de 16,8~V quando carregada e circuito de proteção (BMS). Essa
química foi escolhida pela relação entre energia e massa. Uma bateria de
chumbo-ácido de 12~V e 7~Ah também foi avaliada e descartada: ela pesa
cerca de 2,1~kg, contra 0,3~kg a 0,5~kg da LiPo, e sua tensão de 11,5~V a
12,8~V é baixa demais para alimentar o conversor de 12~V.

A tensão da bateria é distribuída por dois conversores CC-CC abaixadores
(\textit{buck}), cada um com uma chave própria:

\begin{itemize}
	\item um conversor de 12~V alimenta os dois MD10C (chave SW1). Como a
	tensão nos drivers fica regulada na tensão nominal dos motores, o
	ciclo de trabalho pode chegar a 100\% sem sobretensão, e a velocidade
	para um mesmo ciclo de trabalho deixa de depender da carga da
	bateria;
	\item um conversor de 5~V e 5~A alimenta o Raspberry Pi (chave SW2,
	com voltímetro), que por sua vez alimenta pelas portas USB o LiDAR, a
	câmera e o adaptador Wi-Fi.
\end{itemize}

Separar a alimentação dos motores da alimentação da lógica evita que os
picos de corrente na partida dos motores derrubem a tensão do computador.
Mesmo assim, durante os testes o Raspberry Pi indicou subtensão
(\texttt{vcgencmd get\_throttled} retornou \texttt{0x70005}). Isso
provocou redução de frequência do processador e reconexões dos
dispositivos USB. A causa provável é a queda de tensão no cabo micro-USB
e no fusível de entrada da placa quando o LiDAR, a câmera e o Wi-Fi
funcionam juntos. A correção prevista é alimentar a placa diretamente
pelos pinos de 5~V da barra de GPIO, com fio curto e de bitola maior, e,
se necessário, trocar o conversor por um ajustável regulado entre 5,15~V
e 5,2~V. % TODO: atualizar quando a correcao for feita e testada.

\section{Integração dos componentes}
\label{sec:integracao-hw}

% Quando a figura abaixo for descomentada, troque o inicio da frase por
% "A Figura~\ref{fig:pinagem} e a Tabela~\ref{tab:pinagem} mostram".
A Tabela~\ref{tab:pinagem} mostra como os
componentes se ligam ao Raspberry Pi. Eles usam três tipos de interface:

\begin{itemize}
	\item \textbf{GPIO}: os sinais PWM e DIR dos drivers e os canais A e B
	dos encoders. Os encoders são alimentados com 3,3~V, então seus sinais
	já chegam no nível lógico do Raspberry Pi, sem conversor de nível;
	\item \textbf{UART}: a IMU BNO055;
	\item \textbf{USB}: o LiDAR, pelo adaptador CP210x, a câmera Astra Pro
	e o adaptador Wi-Fi.
\end{itemize}

\begin{table}[htb]
	\centering
	\caption{Ligações no Raspberry Pi 3 (numeração BCM e pino físico)}
	\label{tab:pinagem}
	\begin{tabular}{|l|c|c|}
		\hline
		\textbf{Função} & \textbf{GPIO} & \textbf{Pino} \\ \hline
		MD10C esquerdo -- PWM & 12 & 32 \\ \hline
		MD10C esquerdo -- DIR & 17 & 11 \\ \hline
		MD10C direito -- PWM & 13 & 33 \\ \hline
		MD10C direito -- DIR & 22 & 15 \\ \hline
		Encoder dianteiro esquerdo -- A / B & 23 / 24 & 16 / 18 \\ \hline
		Encoder dianteiro direito -- A / B & 25 / 5 & 22 / 29 \\ \hline
		BNO055 -- TX / RX & 14 / 15 & 8 / 10 \\ \hline
		Buzzer & 18 & 12 \\ \hline
		Alimentação dos encoders e da IMU & 3,3~V & 1 / 17 \\ \hline
	\end{tabular}
	\legend{Fonte: Autor.}
\end{table}

% Copie pinagem_robo4wd.png para a pasta figuras/ e descomente:
% \begin{figure}[htb]
% 	\centering
% 	\includegraphics[width=\textwidth]{figuras/pinagem_robo4wd.png}
% 	\caption{Pinagem do Raspberry Pi 3 no robô}
% 	\label{fig:pinagem}
% 	\legend{Fonte: Autor.}
% \end{figure}

Todos os circuitos compartilham o mesmo terra, que é a referência comum
para os sinais lógicos entre o Raspberry Pi, os drivers e os sensores. Os
pinos de terra foram distribuídos pela barra apenas para deixar cada fio
de terra próximo do seu sinal. Os fios de potência dos motores ficam
separados dos fios dos encoders para reduzir o ruído induzido nos sinais
de quadratura.

Na escolha dos pinos, dois pontos foram levados em conta. Primeiro, o
GPIO18, usado pelo buzzer, está ligado internamente ao mesmo canal de PWM
(PWM0) que o GPIO12. Por isso o buzzer é acionado apenas em liga/desliga
ou por PWM por software, por meio de um transistor BC548 com resistor de
base de 1~k$\Omega$. Segundo, a saída de áudio analógico do Raspberry Pi
usa os mesmos periféricos de PWM e foi desabilitada para não interferir
nos motores. O esquemático completo foi desenhado no Proteus (revisão
1.1).
